\documentclass[11pt, a4paper]{article}
\usepackage[utf8]{inputenc}
\usepackage{graphicx}
\usepackage{hyperref}
\usepackage{geometry}
\usepackage{booktabs}
\usepackage{amsmath}

\geometry{margin=2.5cm}

\title{
    \vspace{2cm}
    \textbf{EE200: Signals, Systems and Networks} \\
    \vspace{0.5cm}
    \large Professor: Tushar Sandhan \\
    Course Project — Question 3 \\
    \vspace{1cm}
    \textbf{\Large Q3: Sonic Signatures \& Zapptain America}
    \vspace{2cm}
}

\author{
    \begin{tabular}{lr}
        \textbf{Name:} Jani Ravi Kailash & \textbf{Roll No.:} 240486 \\
        \textbf{Name:} Shreshthraj Bhidodiya & \textbf{Roll No.:} 240996
    \end{tabular}
}
\date{}

\begin{document}

\maketitle

\begin{center}
    \textbf{Live App:} \url{https://ee200-q3-u6rgv2jyrhfqdr4h4lfvje.streamlit.app} \\
    \textbf{Source Code:} \url{https://github.com/Ravi5422/ee200-q3}
\end{center}

\newpage

\section{Introduction}
For this project, we set out to build our own version of Shazam from scratch. The challenge was pretty straightforward but technically deep: if we take a short 5 to 10-second audio clip recorded in a noisy room, how can we reliably figure out which song it belongs to out of a database of 50 tracks?

Instead of relying on pre-built machine learning models or audio fingerprinting libraries, we built the entire system using fundamental Digital Signal Processing (DSP) techniques. We had to find a way to extract a unique "fingerprint" from the audio that wouldn't get destroyed by background noise or volume changes.

\section{How Our Pipeline Works}
To make this happen, we process the raw audio through a four-stage pipeline. This compresses megabytes of raw audio into just a few kilobytes of highly searchable data.

\subsection{Step 1: The Spectrogram}
First, we convert the audio waveform into a time-frequency representation using the Short-Time Fourier Transform (STFT). We used a 1024-sample Hann window with a 50\% overlap. 
We specifically convert the magnitudes to a logarithmic scale (dB). This is super important because human hearing is logarithmic—taking the log allows us to mathematically "see" the quieter, structural harmonics of the song that would otherwise be drowned out by loud bass frequencies.

\subsection{Step 2: Constellation Extraction}
A spectrogram still has too much data. To make it searchable, we isolate only the most prominent peaks. We slide a 20x20 neighborhood window across the spectrogram and keep a point only if it's the absolute maximum in that window and it exceeds the overall average volume. 
This leaves us with a "constellation map"—a sparse scatter plot of dots. These dots usually represent strong instrument notes or drum hits, which are exactly the parts of a song most likely to survive heavy background noise.

\subsection{Step 3: Combinatorial Hashing}
A single peak isn't unique. To identify a song, we look at the relationship between peaks. For every "anchor" peak, we pair it with a few "target" peaks that happen shortly after it. 
We hash these pairs into a format: $(f_{\text{anchor}}, f_{\text{target}}, \Delta t)$.
Because we use the time difference ($\Delta t$) between the peaks instead of their absolute timestamps, the hash doesn't care when the recording started. And because we use frequency bins rather than amplitude, it completely ignores volume changes.

\subsection{Step 4: Histogram Voting}
When we test a query clip, we generate its hashes and look them up in our database. For every match, we calculate the time offset: $\delta = t_{\text{database}} - t_{\text{query}}$. 
If a song is a true match, a ton of hashes will all point to the exact same starting time offset, creating a massive spike in a histogram. Random coincidences just create flat, scattered noise. The song with the tallest spike is our winner.

\section{Database Stats}
We indexed a total of 50 songs (primarily by The Beatles and Queen). 
\begin{itemize}
    \item \textbf{Total unique paired hashes:} 1,495,644
    \item \textbf{Database file size:} 36.25 MB
    \item \textbf{Indexing time:} ~30 seconds
\end{itemize}

\section{Experiments and Observations}

\subsection{Why a Standard DFT Isn't Enough}
We experimented with taking a standard 1D Fourier Transform of the song, but it completely strips away time. It tells you what frequencies exist, but not \textit{when} they happen. Two completely different songs in the same key can have very similar DFTs. A spectrogram is absolutely necessary because the temporal evolution is what actually makes a song unique.

\subsection{The Time-Frequency Trade-off}
We hit a classic DSP hurdle: the Heisenberg uncertainty principle. If we used a really short STFT window (e.g., 256 samples), we got great time precision for drum beats but terrible pitch resolution. If we used a massive window (4096 samples), the pitch was perfect but fast notes blurred together. We found that 1024 samples gave us the perfect middle ground to find clean constellation peaks.

\subsection{Noise Robustness}
We tested the system by injecting massive amounts of white noise into the query clips. Impressively, it maintained a \textbf{100\% accuracy rate} even at 0 dB SNR (meaning the noise was exactly as loud as the music itself). The system only started dropping matches at -3 dB. This works so well because random noise generates random hashes that fail to form a coherent histogram spike, while the true song peaks consistently vote together.

\subsection{Pitch Shifting vs. Time Stretching}
We also tested how the algorithm handles distorted audio. 
Pitch shifting completely breaks the system. If you shift the pitch, all the frequencies multiply, meaning every single peak falls into a new bin and the hashes no longer match the database. 
Time stretching, on the other hand, is a bit more forgiving. Since the frequencies don't change, only the $\Delta t$ expands. Small stretches ($\pm$10\%) are tolerated because the $\Delta t$ still falls into the same discrete integer bins, but heavy stretching causes it to fail.

\section{Zapptain America (Our Web App)}
We didn't just write scripts; we packaged the entire system into a modern, highly responsive web application using Streamlit. 

We gave the interface a sleek dark-mode aesthetic with custom CSS. One of the coolest features is the Library tab: it dynamically displays all 50 songs, and we generated custom, brightly colored constellation maps to act as the background image for each song card, creating a beautiful mosaic effect.

To make sure it could be reliably deployed to the cloud, we pre-compiled the constellation data into a tiny file. This means the live website doesn't need to process hundreds of megabytes of raw MP3s to draw the visual plots—it just loads the math instantly.

\section{Conclusion}
Building this from scratch was incredibly rewarding. By extracting sparse constellation peaks, encoding their relative relationships, and using a histogram voting system, we managed to achieve Shazam-level accuracy using nothing but basic math arrays and signal processing logic. The final system is lightning-fast, highly resistant to noise, and beautifully packaged for the web.

\end{document}
